\section{Introduction}

Ce TD vous place en situation de bureau d'études : vous devez concevoir les
éléments d'architecture et de fonctionnement d'une cellule robotique de
traitement de pièces, avant sa mise en œuvre effective sur le système réel.

\begin{UPSTIinfor}{Organisation de la séance}
\begin{itemize}
    \item Travail en \textbf{équipe de 3} ;
    \item Durée totale : \textbf{4h} ;
    \item Un \textbf{rapport unique par équipe} est à rendre en fin de séance,
    reprenant l'ensemble des livrables demandés dans ce sujet ;
    \item Le document ressource \emph{EtherNet/IP} ainsi que le cours au tableau
    sur ce protocole constituent les supports théoriques nécessaires à la
    partie réseau de ce TD.
\end{itemize}
\end{UPSTIinfor}

\begin{UPSTIwarning}{Ce qui est attendu de vous}
Ce sujet ne décrit volontairement \textbf{pas} la solution technique. C'est à
vous de concevoir l'architecture, le réseau et les modes de fonctionnement à
partir du besoin exprimé et des contraintes imposées. Les questions posées
sont là pour guider votre démarche de conception, pas pour vous donner la
réponse.
\end{UPSTIwarning}

Une répartition indicative du temps de travail est proposée ci-dessous ; elle
peut être adaptée par votre équipe en fonction de votre avancement.

\begin{table}[h!]
    \centering
    \begin{tabular}{|l|l|}
        \hline
        \textbf{Partie} & \textbf{Durée indicative} \\
        \hline
        Présentation du système et prise en main du besoin & 20 min \\
        \hline
        Schéma synoptique et architecture réseau (livrables 1 et 2) & 1h00 \\
        \hline
        Modes de fonctionnement GMMA (livrable 3) & 1h00 \\
        \hline
        Cartographie mémoire M340 et sécurités (livrables 4 et 5) & 1h00 \\
        \hline
        Trajectoires du robot (livrable 6) & 30 min \\
        \hline
        Mise en forme et relecture du rapport & 10 min \\
        \hline
    \end{tabular}
    \caption{Répartition indicative du temps de travail sur 4h}
\end{table}
